
% 3. Environnement unique et superposition directe
\begin{tikzpicture}
  % On dessine les deux éléments sur le MÊME repère absolu
  \clip[](-1.2, -1) rectangle (1.1 , 1.3) ;

  \node[transform canvas={scale=1}] at (0,0) {
    \begin{tikzpicture}

        \begin{scope}[xscale=1,yscale=1]
            % Grille principale
            \draw[opacity=0.2,very thin,line width=0.01cm,color=colorThree](-0.15,-1.95)grid(1.95,0.15);
            % Sous-grille
            \draw[opacity=0.2,very thin,line width=0.01cm,color=colorThree](-0.15,-2)grid[step=0.1](1.95,0.15);

            % Graduations X
            \foreach \x in {-0.1,0,0.1,...,2} {
                \draw[line width=0.01cm,color=colorOne](\x,0-0.03)--(\x,0+0.03);
            }
            \foreach \x in {1,...,1} {
                \draw[line width=0.01cm,color=colorOne](\x,0-0.05)--(\x,0+0.05)
                node[shift={(0,0.05)},font=\fontsize{0.05cm}{0.05cm}\bfseries\selectfont]{\x};
            }
            % Graduations Y
             \foreach \y in {0.1,0,-0.1,...,-2} {
                \draw[line width=0.01cm,color=colorOne,](0-0.03,\y)--(0+0.03,\y);
            }
            \foreach \y in {-1,...,-1} {
                \draw[line width=0.01cm,color=colorOne,](0-0.05,\y)--(0+0.05,\y)
                node[shift={(-0.15,0)},font=\fontsize{0.05cm}{0.05cm}\bfseries\selectfont]{\y};
            }

        \end{scope}
        
        % 2. Axe des abscisses (x)
        \draw[line width=0.01cm,color=colorOne,->,
        >={Computer Modern Rightarrow[length=0.04cm, width=0.06cm]}
        ] (-0.2,0) -- (2,0) 
            node[pos=1.0, xshift=-0.06cm,yshift=-0.05cm, color=colorOne, font=\fontsize{0.05cm}{0.05cm}\bfseries\selectfont] {$x$};

        % 3. Axe des ordonnées (S(x))
        \draw[line width=0.01cm,color=colorOne,->,
        >={Computer Modern Rightarrow[length=0.04cm, width=0.06cm]}
        ] (0,-2) -- (0,0.2) 
            node[pos=1.0, xshift=-0.05cm,yshift=-0.06cm, color=colorOne, font=\fontsize{0.05cm}{0.05cm}\bfseries\selectfont] {$y$}
            ;

        % 4. Courbe de la parabole y = x^2 (domaine restreint pour rester dans le cadre)
        %\draw[smooth, line width=1.0ex, color=colorFour, domain=-2.1:2.1, samples=100] plot (\x, {1-2*\x*\x});

        \def\vzero{3} \def\g{9.81} \def\h{0} \def\dt{0.05} 

        % \def\listeLabels{
        %     0//O/,
        %     1//M_1///,
        %     2//M_2/\overrightarrow{OM}_2//,
        %     3/////,
        %     4/////,
        %     5//M_5/\overrightarrow{OM}_5//,
        %     6//M_6///,
        %     7//M_7/\overrightarrow{OM}_7///,
        %     8//M_8//\vec{v}_8/,
        %     9//M_9/\overrightarrow{OM}_9//\vec{a}_9/,
        %     10//M_{10}//\vec{v}_{10}/,
        %     11//M_{11}/\overrightarrow{OM}_{11}///,
        %     12//M_{12}////,
        %     13//M_{13}////
        % }

        \def\listeLabels{
            0//O///,
            1//M_1///,
            2//M_2//,
            3/////,
            4/////,
            5//M_5/\overrightarrow{OM}_5//,
            6//M_6//\vec{v}(t_6)/,
            7//M_7/\overrightarrow{OM}_7//,
            8//M_8///,
            9//M_9///,
            10//M_{10}///,
            11//M_{11}///,
            12//M_{12}///,
        }

        \foreach \n/\labelBottom/\labelTop/\LabelvecPos/\LabelvecVit/\LabelvecAcc in \listeLabels {%

            \pgfmathsetmacro{\t}{\n*\dt}
            \pgfmathsetmacro{\x}{\vzero*\t}
            \pgfmathsetmacro{\y}{\h - 0.5*\g*\t*\t}

            % Point
            \fill[color=colorFour] (\x,\y) circle (0.01);

            % Label inférieur
            \if\relax\detokenize{\labelBottom}\relax
            \else
                \node[
                    below,
                    font=\fontsize{0.05cm}{0.05cm}\bfseries\selectfont
                ] at (\x,\y) {$\labelBottom$};
            \fi

            % Label supérieur
            \if\relax\detokenize{\labelTop}\relax
            \else
                \node[
                    font=\fontsize{0.05cm}{0.05cm}\bfseries\selectfont
                ] at (\x+0.05,\y+0.05) {$\labelTop$};
            \fi

            % Vecteur position
            \if\relax\detokenize\expandafter{\LabelvecPos}\relax
            \else
                \draw[opacity = 0.3, line width=0.01cm,color=colorFour,->,>={Computer Modern Rightarrow[length=0.04cm,width=0.06cm]}](0,0) -- (\x,\y)
                node[pos=0.7,xshift=0cm,yshift=0.05cm,color=colorFour,font=\fontsize{0.05cm}{0.05cm}\bfseries\selectfont] {$\LabelvecPos$};
            \fi

            % Vecteur vitesse
            \if\relax\detokenize\expandafter{\LabelvecVit}\relax
            \else
                \draw[shift = {(\x,\y)}, line width=0.01cm,color=colorFive,->,>={Computer Modern Rightarrow[length=0.04cm,width=0.06cm]}](0,0) -- (\vzero/10,-\g*\t/10)
                node[pos=0.5,xshift=0cm,yshift=0.1cm,color=colorFive,font=\fontsize{0.05cm}{0.05cm}\bfseries\selectfont] {$\LabelvecVit$};
            \fi

            % Vecteur accélération
            \if\relax\detokenize\expandafter{\LabelvecAcc}\relax
            \else
                \draw[shift = {(\x,\y)},line width=0.01cm,color=colorSix,->,>={Computer Modern Rightarrow[length=0.04cm,width=0.06cm]}](0,0) -- (0,-\g/10)
                node[pos=0.5,xshift=0cm,yshift=0.1cm,color=colorSix,font=\fontsize{0.05cm}{0.05cm}\bfseries\selectfont] {$\LabelvecAcc$};
            \fi
        }
        
        %echelle
        \draw[line width=0.01cm,shift={(1.7,-0.3)}]
            (-0.1,0)edge[shift={(0,0)},color=colorFour,{Latex[length=0.4mm,width=0.3mm]}-{Latex[length=0.4mm,width=0.3mm]}]node[pos=0.5,xshift=0cm,yshift=0.03cm,color=colorFour,font=\fontsize{0.03cm}{0.02cm}\bfseries\selectfont]{$20~cm$}(0.1,0)
            (-0.1,0)edge[shift={(0,-0.1)},color=colorFive,{Latex[length=0.4mm,width=0.3mm]}-{Latex[length=0.4mm,width=0.3mm]}]node[pos=0.5,xshift=0cm,yshift=0.03cm,color=colorFive,font=\fontsize{0.03cm}{0.01cm}\bfseries\selectfont]{$2~m \!\cdot \! s^{-1}$}(0.1,0)
            (-0.1,0)edge[shift={(0,-0.2)},color=colorSix,{Latex[length=0.4mm,width=0.3mm]}-{Latex[length=0.4mm,width=0.3mm]}]node[pos=0.5,xshift=0cm,yshift=0.03cm,color=colorSix,font=\fontsize{0.03cm}{0.02cm}\bfseries\selectfont]{$2~m\!\cdot \! s^{-2}$}(0.1,0)
        ;


        \draw[opacity = 0.5,line width=0.01cm,shift={(0,0)}]
            (\vzero*\dt*5,\h - 0.5*\g*\dt*\dt*5*5)edge[shift={(0,0)},color=colorFour,-{Computer Modern Rightarrow[length=0.04cm,width=0.06cm]}]node[pos=0.5,above=-1.4mm,color=colorFour,font=\fontsize{0.05cm}{0.04cm}\bfseries\selectfont, rotate = -45]{$ \Delta\overrightarrow{OM}{(t_{5\to 7})}$}(\vzero*\dt*7,\h - 0.5*\g*\dt*\dt*7*7)
        ;



        




        
        
        %\foreach \n in {0,...,7} { \pgfmathsetmacro{\t}{\n*\dt} \pgfmathsetmacro{\x}{\vzero*\t} \pgfmathsetmacro{\y}{\h - 0.5*\g*\t*\t} 
        %\fill[color=colorFour] (\x,\y) circle (0.8); }

        % \draw[ opacity = 0.5, dashed, smooth, line width=1.0ex, color=colorFour, domain=0:1, samples=100 ] plot ({\vzero*\x}, {\h - 0.5*\g*\x*\x});

        % \draw[ dashed, smooth, line width=1.0ex, color=colorFour, domain=0:1, samples=100 ] plot (\vzero, {-\g*\x});

        % \node[below right, color=colorFour, font=\fontsize{15pt}{10pt}\bfseries\selectfont , rotate = 80] at (-2.4, -4.4) {$z(t) = 1-2t^2$};

        % \draw[dashed, smooth, line width=0.8ex, color=colorFive, domain=-0.5:3, samples=100] 
        %     plot (\x, {-4*\x+3});

        % \node[below right, color=colorFive, font=\fontsize{7pt}{10pt}\bfseries\selectfont , rotate = -75] at (2.5, -5.0) {$\dot{z}(a)(t-a)+z(a)$};

        % \draw[line width=0.8ex, color=colorOne] (1,0.3) -- (1,-0.3) 
        %     node[pos=-0.1, above , color=colorOne, font=\fontsize{20pt}{10pt}\bfseries\selectfont] {$a$};
        % \draw[line width=0.8ex, color=colorOne] (2,0.3) -- (2,-0.3) 
        %     node[pos=-0.1, above, color=colorOne, shift={(0.8,-0.0)},  font=\fontsize{20pt}{10pt}\bfseries\selectfont] {$a + \mathrm{d} a$};

        % \draw[dashed, line width=0.5ex, color=colorOne] (1,0) -- (1,-8);
        % \draw[dashed, line width=0.5ex, color=colorOne] (-2.5,-1) -- (4,-1)
        %     node[pos=-0.0, left, color=colorOne, font=\fontsize{20pt}{10pt}\bfseries\selectfont] {$z(a)$};
        % \draw[dashed, line width=0.5ex, color=colorOne] (2,0) -- (2,-8);
        % \draw[dashed, line width=0.5ex, color=colorOne] (-2.5,-7) -- (4,-7)
        %     node[pos=-0.0, left, color=colorOne, font=\fontsize{20pt}{10pt}\bfseries\selectfont] {$z(a + \mathrm{d} a )$};

        % \draw[dashed, line width=0.5ex, color=colorOne] (-2.5,-5) -- (4,-5)
        %     node[pos=-0.0, left, color=colorOne, font=\fontsize{20pt}{10pt}\bfseries\selectfont] {$\dot{z}(a)\mathrm{d} a+z(a)$};

        % \draw[dashed, line width=0.5ex, color=colorOne , -{Latex[length=4mm,width=3mm]}] (1,1.2) -- (1.2,1.2);
        % \draw[dashed, line width=0.5ex, color=colorOne , {Latex[length=4mm,width=3mm]}- ] (1.8,1.2) -- (2,1.2);
        % \node[color=colorOne, font=\fontsize{10pt}{10pt}\bfseries\selectfont ] at (1.5, 1.2) {$\mathrm{d} a$};

        % \draw[line width=0.5ex, color=colorOne , {Latex[length=4mm,width=3mm]}- ] (3,-5) -- (3,-3.5);
        % \draw[line width=0.5ex, color=colorOne , -{Latex[length=4mm,width=3mm]} ] (3,-2.5) -- (3,-1);
        % \node[color=colorOne, font=\fontsize{10pt}{10pt}\bfseries\selectfont , rotate = 90 ] at (3, -3) {$\mathrm{d} z(a)$};

        % \draw[line width=0.5ex, color=colorOne , {Latex[length=4mm,width=3mm]}- ] (4.0,-7) -- (4.0,-4.5);
        % \draw[line width=0.5ex, color=colorOne , -{Latex[length=4mm,width=3mm]} ] (4,-3.5) -- (4,-1);
        % \node[color=colorOne, font=\fontsize{10pt}{10pt}\bfseries\selectfont , rotate = 90 ] at (4.0,-4) {$\Delta z(a)$};



     \end{tikzpicture}
};
        

  % Grille de repérage
  %\draw[line width=0.5ex , color = red] (-1.2, -1) rectangle (1.1 , 1.3) ;
  %\drawgrid{-6.}{6}{-5.5}{4.5}{1}{1}
\end{tikzpicture}